% GENERATED FILE. Edit canonical lessons, then run npm run generate:books.
% canonical-source-hash: fnv1a64:b59510b317e1dedc
% canonical-lessons: ML-C439-uddharikkuka, ML-C439-samgrahikkuka, ML-C439-upasamharikkuka, ML-C439-avalambam, ML-C439-granthakaran

\chapter{Quoting and Summing Up}
\label{ch:ml-quoting-and-summing-up}
\hlchaptermodality{439}

\begin{chapteropening}
\textbf{By the end of this chapter:} \emph{I can say to quote, to summarise, to conclude, a reference and an author.}

Complete the last lesson of chapter 439: I can say to quote, to summarise, to conclude, a reference and an author.
\end{chapteropening}

\section[\texorpdfstring{\ml{ഉദ്ധരിക്കുക} (uddharikkuka)}{ഉദ്ധരിക്കുക (uddharikkuka)}]{\ml{ഉദ്ധരിക്കുക} (uddharikkuka) — to quote; to cite}

\label{lesson:ML-C439-uddharikkuka}



\begin{quote}
Before the new one: say the Malayalam for to interpret, then the Malayalam for to define.
\end{quote}

\subsection*{You'll want to know: \ml{ഉദ്ധരിക്കുക}}
\textbf{\ml{ഉദ്ധരിക്കുക}} — \emph{uddharikkuka} — \textquotedblleft{}to quote; to cite\textquotedblright{}.

\begin{quote}
\ml{പ്രഭാഷകൻ} \ml{ഒരു} \ml{കവിയുടെ} \ml{വരികൾ} \ml{ഉദ്ധരിച്ചു}.
\end{quote}

\begin{quote}
\emph{prabhāṣakan oru kaviyuṭe varikaḷ uddhariccu.}
\end{quote}

\begin{quote}
The speaker quoted a poet's lines.
\end{quote}

Past: \ml{ഉദ്ധരിച്ചു}, \emph{uddhariccu}.

The noun is \ml{ഉദ്ധരണി}, a quotation.

\begin{cousinweb}
\textbf{\ml{ഉദ്ധരിക്കുക}} (\emph{uddharikkuka}) is built from Sanskrit \emph{uddharaṇa}, \textquotedblleft{}taking up, drawing out\textquotedblright{} (Monier-Williams), with the Malayalam verb ending -\ml{ിക്കുക}: to quote is to draw words out of a book.
\end{cousinweb}

\begin{grammarlens}[title={What you've built}]
One more word for giving reasons and opinions.
\end{grammarlens}

\subsection*{Your turn}
\begin{itemize}
\raggedright
  \item \emph{Say it:} \emph{uddharikkuka}
  \item \emph{Say it:} \emph{uddharikkuka}, once more
  \item \emph{Say it:} \emph{prabhāṣakan oru kaviyuṭe varikaḷ uddhariccu.} — The speaker quoted a poet's lines.
  \item \emph{Recall:} say the Malayalam for to interpret, then the Malayalam for to define, then say \emph{uddharikkuka} again
\end{itemize}

\begin{tcolorbox}[breakable,colback=teal!4,colframe=teal!35!black,title={Before you move on}]
What does \ml{ഉദ്ധരിക്കുക} mean? (\textquotedblleft{}To quote; to cite\textquotedblright{}.) Say it once more.
\end{tcolorbox}
\section[\texorpdfstring{\ml{സംഗ്രഹിക്കുക} (saṁgrahikkuka)}{സംഗ്രഹിക്കുക (saṁgrahikkuka)}]{\ml{സംഗ്രഹിക്കുക} (saṁgrahikkuka) — to summarise}

\label{lesson:ML-C439-samgrahikkuka}



\begin{quote}
Before the new one: say the Malayalam for to define, then the Malayalam for to quote.
\end{quote}

\subsection*{You'll want to know: \ml{സംഗ്രഹിക്കുക}}
\textbf{\ml{സംഗ്രഹിക്കുക}} — \emph{saṁgrahikkuka} — \textquotedblleft{}to summarise\textquotedblright{}.

\begin{quote}
\ml{അവൾ} \ml{ആ} \ml{ലേഖനം} \ml{അഞ്ച്} \ml{വാചകങ്ങളിൽ} \ml{സംഗ്രഹിച്ചു}.
\end{quote}

\begin{quote}
\emph{avaḷ ā lēkhanaṁ añcŭ vācakaṅṅaḷil saṁgrahiccu.}
\end{quote}

\begin{quote}
She summarised that article in five sentences.
\end{quote}

Past: \ml{സംഗ്രഹിച്ചു}, \emph{saṁgrahiccu}.

The noun is \ml{സംഗ്രഹം}, a summary.

\begin{cousinweb}
\textbf{\ml{സംഗ്രഹിക്കുക}} (\emph{saṁgrahikkuka}) is built from Sanskrit \emph{saṁgraha}, \textquotedblleft{}a collection, a holding together\textquotedblright{} (CDIAL 12852), with the Malayalam verb ending -\ml{ിക്കുക}: to gather a text into a little.
\end{cousinweb}

\begin{grammarlens}[title={What you've built}]
One more word for giving reasons and opinions.
\end{grammarlens}

\subsection*{Your turn}
\begin{itemize}
\raggedright
  \item \emph{Say it:} \emph{saṁgrahikkuka}
  \item \emph{Say it:} \emph{saṁgrahikkuka}, once more
  \item \emph{Say it:} \emph{avaḷ ā lēkhanaṁ añcŭ vācakaṅṅaḷil saṁgrahiccu.} — She summarised that article in five sentences.
  \item \emph{Recall:} say the Malayalam for to define, then the Malayalam for to quote, then say \emph{saṁgrahikkuka} again
\end{itemize}

\begin{tcolorbox}[breakable,colback=teal!4,colframe=teal!35!black,title={Before you move on}]
What does \ml{സംഗ്രഹിക്കുക} mean? (\textquotedblleft{}To summarise\textquotedblright{}.) Say it once more.
\end{tcolorbox}
\section[\texorpdfstring{\ml{ഉപസംഹരിക്കുക} (upasaṁharikkuka)}{ഉപസംഹരിക്കുക (upasaṁharikkuka)}]{\ml{ഉപസംഹരിക്കുക} (upasaṁharikkuka) — to conclude; to wind up}

\label{lesson:ML-C439-upasamharikkuka}



\begin{quote}
Before the new one: say the Malayalam for to quote, then the Malayalam for to summarise.
\end{quote}

\subsection*{You'll want to know: \ml{ഉപസംഹരിക്കുക}}
\textbf{\ml{ഉപസംഹരിക്കുക}} — \emph{upasaṁharikkuka} — \textquotedblleft{}to conclude; to wind up\textquotedblright{}.

\begin{quote}
\ml{ഒരു} \ml{ചെറിയ} \ml{കഥയോടെ} \ml{അദ്ദേഹം} \ml{പ്രസംഗം} \ml{ഉപസംഹരിച്ചു}.
\end{quote}

\begin{quote}
\emph{oru ceṟiya kathayōṭe addēhaṁ prasaṁgaṁ upasaṁhariccu.}
\end{quote}

\begin{quote}
He concluded the speech with a short story.
\end{quote}

Past: \ml{ഉപസംഹരിച്ചു}, \emph{upasaṁhariccu}.

The noun is \ml{ഉപസംഹാരം}, the conclusion of a speech or essay.

\begin{cousinweb}
\textbf{\ml{ഉപസംഹരിക്കുക}} (\emph{upasaṁharikkuka}) is built from Sanskrit \emph{upasaṁhāra}, \textquotedblleft{}the act of withdrawing, drawing together\textquotedblright{} (Monier-Williams), with the Malayalam verb ending -\ml{ിക്കുക}.
\end{cousinweb}

\begin{grammarlens}[title={What you've built}]
One more word for giving reasons and opinions.
\end{grammarlens}

\subsection*{Your turn}
\begin{itemize}
\raggedright
  \item \emph{Say it:} \emph{upasaṁharikkuka}
  \item \emph{Say it:} \emph{upasaṁharikkuka}, once more
  \item \emph{Say it:} \emph{oru ceṟiya kathayōṭe addēhaṁ prasaṁgaṁ upasaṁhariccu.} — He concluded the speech with a short story.
  \item \emph{Recall:} say the Malayalam for to quote, then the Malayalam for to summarise, then say \emph{upasaṁharikkuka} again
\end{itemize}

\begin{tcolorbox}[breakable,colback=teal!4,colframe=teal!35!black,title={Before you move on}]
What does \ml{ഉപസംഹരിക്കുക} mean? (\textquotedblleft{}To conclude; to wind up\textquotedblright{}.) Say it once more.
\end{tcolorbox}
\section[\texorpdfstring{\ml{അവലംബം} (avalambaṁ)}{അവലംബം (avalambaṁ)}]{\ml{അവലംബം} (avalambaṁ) — a reference; a support}

\label{lesson:ML-C439-avalambam}



\begin{quote}
Before the new one: say the Malayalam for to summarise, then the Malayalam for to conclude.
\end{quote}

\subsection*{You'll want to know: \ml{അവലംബം}}
\textbf{\ml{അവലംബം}} — \emph{avalambaṁ} — \textquotedblleft{}a reference; a support\textquotedblright{}.

\begin{quote}
\ml{ലേഖനത്തിന്റെ} \ml{അവസാനം} \ml{അവലംബം} \ml{കൊടുത്തിട്ടുണ്ട്}.
\end{quote}

\begin{quote}
\emph{lēkhanattinṟe avasānaṁ avalambaṁ koṭuttiṭṭuṇṭŭ.}
\end{quote}

\begin{quote}
The references are given at the end of the article.
\end{quote}

It is first a support, something to lean on.

\begin{cousinweb}
\textbf{\ml{അവലംബം}} (\emph{avalambaṁ}) is borrowed from Sanskrit \emph{avalamba}, \textquotedblleft{}hanging down; a stay, a support\textquotedblright{} (CDIAL 826). A reference is what a writer leans on.
\end{cousinweb}

\begin{grammarlens}[title={What you've built}]
One more word for giving reasons and opinions.
\end{grammarlens}

\subsection*{Your turn}
\begin{itemize}
\raggedright
  \item \emph{Say it:} \emph{avalambaṁ}
  \item \emph{Say it:} \emph{avalambaṁ}, once more
  \item \emph{Say it:} \emph{lēkhanattinṟe avasānaṁ avalambaṁ koṭuttiṭṭuṇṭŭ.} — The references are given at the end of the article.
  \item \emph{Recall:} say the Malayalam for to summarise, then the Malayalam for to conclude, then say \emph{avalambaṁ} again
\end{itemize}

\begin{tcolorbox}[breakable,colback=teal!4,colframe=teal!35!black,title={Before you move on}]
What does \ml{അവലംബം} mean? (\textquotedblleft{}A reference; a support\textquotedblright{}.) Say it once more.
\end{tcolorbox}
\section[\texorpdfstring{\ml{ഗ്രന്ഥകാരൻ} (granthakāran)}{ഗ്രന്ഥകാരൻ (granthakāran)}]{\ml{ഗ്രന്ഥകാരൻ} (granthakāran) — an author}

\label{lesson:ML-C439-granthakaran}



\begin{quote}
Before the new one: say the Malayalam for to conclude, then the Malayalam for a reference.
\end{quote}

\subsection*{You'll want to know: \ml{ഗ്രന്ഥകാരൻ}}
\textbf{\ml{ഗ്രന്ഥകാരൻ}} — \emph{granthakāran} — \textquotedblleft{}an author\textquotedblright{}.

\begin{quote}
\ml{ഗ്രന്ഥകാരൻ} \ml{പുസ്തകത്തിൽ} \ml{ഒപ്പിട്ടു} \ml{തന്നു}.
\end{quote}

\begin{quote}
\emph{granthakāran pustakattil oppiṭṭu tannu.}
\end{quote}

\begin{quote}
The author signed the book for me.
\end{quote}

A woman author is \ml{ഗ്രന്ഥകാരി}.

\begin{cousinweb}
\textbf{\ml{ഗ്രന്ഥകാരൻ}} (\emph{granthakāran}) is built from Sanskrit \emph{grantha}, \textquotedblleft{}a knot; a book\textquotedblright{} (CDIAL 4350), and \emph{kāra}, \textquotedblleft{}making, doing\textquotedblright{} (CDIAL 3053), with the Malayalam masculine ending -\ml{ൻ}: the maker of a book.
\end{cousinweb}

\begin{grammarlens}[title={What you've built}]
One more word for giving reasons and opinions.
\end{grammarlens}

\subsection*{Your turn}
\begin{itemize}
\raggedright
  \item \emph{Say it:} \emph{granthakāran}
  \item \emph{Say it:} \emph{granthakāran}, once more
  \item \emph{Say it:} \emph{granthakāran pustakattil oppiṭṭu tannu.} — The author signed the book for me.
  \item \emph{Recall:} say the Malayalam for to conclude, then the Malayalam for a reference, then say \emph{granthakāran} again
\end{itemize}

\begin{tcolorbox}[breakable,colback=teal!4,colframe=teal!35!black,title={Before you move on}]
What does \ml{ഗ്രന്ഥകാരൻ} mean? (\textquotedblleft{}An author\textquotedblright{}.) Say it once more.
\end{tcolorbox}
